\section{基本FT设计}

\insertfigureone

Figure 1展示了我们用于fault-tolerant VM的系统基本设置。%
对于一个希望提供fault tolerance能力的VM(被称为primary VM)，%
我们会在另一台物理服务器上运行一个backup VM，使它保持同步，%
并以与primary VM完全相同的方式执行，只是有很小的时间滞后。%
我们称这两个VM处于virtual lockstep\emph{(译者认为可以称为%
“步调一致”)}状态。VM的virtual disk位于%
shared storage上(例如Fibre Channel或iSCSI disk %
array)，因此primary VM和backup VM都可以访问这些磁盘以进行输%
入和输出(我们将在section 4.1讨论一种primary VM和backup VM使用独%
立的非共享virtual disk的设计)。只有primary VM会在网络上宣告自%
己的存在，因此所有network input\emph{(译者注：网络报文)}都%
会到达primary VM。类似地，所有其他输入(例如键盘和鼠标)也只%
会到达primary VM。\emph{(译者注：如果通过VNC连接VM内的Linux%
那么键盘和鼠标也算network input；当然vSphere Web Console/VMRC的控制台%
会直接向VM的虚拟键盘/鼠标设备注入事件，因此确实存在%
keyboard/mouse input。VMware官方资料也明确讨论了Web Console%
和VMRC中的键鼠事件。读者可以暂时不感知keyboard/mouse，就假设咱们%
只能VNC，通篇论文只关注网络和磁盘相关即可)}%

primary VM收到的所有输入都会通过一条称为\textbf{logging channel}的%
网络连接发送给backup VM。%
对于服务器的工作负载而言，主要的输入流量来自网络和磁盘。如下文section 2.1所述，%
为了确保backup VM与primary VM以相同的方式执行%
non-deterministic operation，必要时还会传输额外信息
\emph{(译者注：这里的关键是，backup VM%
不能让自己所在物理机上的外部事件直接决定guest-visible execution。%
比如backup host自己产生的物理timer interrupt、自己收到的外部网络包，%
都不能直接作为backup guest OS的输入。对于guest-visible的timer interrupt、%
I/O completion interrupt、网络输入、磁盘读结果等，primary会记录相关%
log entry，并通过logging channel发送给backup。backup在replay时按照%
log entry指定的内容和执行位置插入这些输入或异步事件，从而消除状态偏移。%
在默认设计中，disk read结果也作为输入由primary通过logging channel发送给backup；%
论文后文还讨论了一种让backup自己执行disk read的替代设计)}。%
结果是，backup VM始终以与primary VM完全相同的方式执行。%

不过，backup VM的输出会被hypervisor丢弃，%
因此只有primary VM会产生实际输出并返回给客户端。如section 2.2所述，%
primary VM和backup VM遵循一个特定协议，其中包括backup VM的%
显式acknowledgment，以确保primary VM failure时不会丢%
失数据。%
\\\emph{译者注：“丢数据”指的是primary已经对外产生了某个可见结果，%
但backup没有收到足够的log来重放到这个结果对应的状态。}%

为了检测primary VM或backup VM是否已经发生故障，我们的%
系统结合使用了两种机制：相关服务器之间的心跳检测，以及%
对logging channel上流量的监控。除此之外，即使出现%
primary服务器和backup彼此失去通信的这种split-brain(脑裂)情况%
\emph{(这里的split-brain不是指系统真的产生了两个合法primary，%
而是指primary和backup所在服务器之间失去通信后，backup可能误判%
primary已经故障并尝试go live；但primary可能其实仍在运行并继续%
对外服务。如果此时backup也开始对外输出，就会出现两个VM同时对外服%
务，导致网络响应、磁盘写入和客户端状态混乱)}，%
我们也必须确保primary VM和backup VM中只有一个会接管执行%
\emph{(译者注：后文会说明，这里利用了共享磁盘上的test-and-set操作)}。%

在下面各节中，我们将对几个重要方面给出更多细节。在section 2.1，%
我们给出deterministic replay技术的一些细节，%
该技术通过logging channel发送的信息确保primary VM和%
backup VM保持同步。在section 2.2，我们描述FT协议的一条基本规则，%
它确保primary VM failure时不会丢失数据。在section 2.3，%
我们描述用于以正确方式检测并响应failure的方法。

\subsection{deterministic replay的实现}

如前所述，复制服务器(或VM)的执行可以被建模为复制一个%
deterministic state-machine。如果两个%
deterministic state-machine从相同的初始状态开始，%
并以相同顺序获得完全相同的输入，那么它们将经历相同的状态序列，并产生相同的输出。%
VM具有广泛的input集合，包括传入的network packet、disk read，
以及来自键盘和鼠标的input。%
non-deterministic event(例如virtual interrupt)%
和non-deterministic operation(例如%
读取处理器的时钟周期计数器)也会影响VM的状态。%
由于VM可能运行任意操作系统和任意工作负载\emph{(译者注：不能假定VM只运行%
类似“纯计算”的任务，也得考虑是否会读取系统时间，或者进行某些依赖系统%
时间片的工作。实际上network就是这样的工作，TCP协议其实内部会依赖定时器)}，%
这为VM的执行带来了三个挑战：(1)正确捕获保证%
backup VM deterministic execution所需的所有输入和%
non-determinism；(2)正确地把这些输入和%
non-determinism应用到backup VM；(3)%
以不降低性能的方式做到这一点。此外，x86微处理器中的许多复杂操作%
具有未定义的、因而也是non-deterministic的副作用。%
捕获这些未定义副作用并replay它们以产生相同状态，是另一个挑战。%
\\\emph{%
译者注：这里的disk read是从外部磁盘向VM引入数据的操作。guest OS driver%
通常会在guest memory中的descriptor中描述请求，再通过MMIO doorbell通知%
虚拟设备；具体队列布局因设备模型而异。对于read，设备或hypervisor取得磁盘%
数据后，会在确定的completion投递点更新guest memory和虚拟设备状态，再向%
guest注入completion interrupt。VMware FT需要记录read数据、completion结果%
及其精确投递位置，使backup在相同执行点看到相同结果。guest在completion之前%
访问目标内存仍可能看到旧数据，这正是section 3.4使用bounce buffer隔离异步%
DMA更新的原因。对于write，数据从guest memory流向磁盘，属于external output；%
其completion状态和投递时机仍可能具有non-determinism，也需要记录必要信息。}%

VMware deterministic replay[15]正是为VMware %
vSphere平台上的x86 VM提供了这种功能。%
deterministic replay会把VM的输入以及与VM执行相关的所有可能%
的non-determinism，record为一串log entry并写入log %
file。之后通过读取文件中的log entry，可以精确replay VM的执行。%
对于non-deterministic operation，%
会记录足够的信息，使该操作能够以相同的状态变化和输出被复现。%
对于timer interrupt或IO完成interrupt等%
non-determinism事件，还会record事件发生时的精确指令。%
在replay期间，该事件会在指令流中的同一位置被投递。%
VMware deterministic replay实现了一种高效的事件%
record和事件投递机制，采用了多种技术，包括与AMD[2]和Intel[8]%
共同开发的硬件性能计数器\emph{(译者注：硬件性能计数器是用%
来分析性能的，它里面有很多事件可以使用，但要实现VM FT只需关注%
instructions retired就行，它代表已提交的机器指令数)}。%

Bressoud和Schneider[3]提到把VM的执行划分为若干epoch，%
诸如interrupt这样的non-deterministic event只在一个%
epoch结束时投递。epoch的概念似乎被用作一种批处理机制，%
因为在事件发生的精确指令处分别投递每个interrupt成本太高。然而，%
我们的事件投递机制足够高效，因此VMware deterministic replay%
不需要使用epoch。每个interrupt都会在发生时被记录，%
并在replay时高效地投递到适当的指令处。
\\\emph{译者注：先简要描述Bressoud和Schneider的epoch机制。%
该系统同样采用primary/backup架构，两台机器上都有一种特殊寄存器，%
该寄存器会在每执行完一条指令后递减一次；当其变为负数时，%
处理器会触发一个控制中断，使hypervisor重新获得控制权。%
例如，如果将该寄存器设置为1000，那么按照字面语义，VM会在%
执行完1001条指令后结束当前epoch，并陷入由该特殊寄存器触发的%
控制中断(hypervisor层的中断)。此时hypervisor可以统一处理本epoch内缓存起来的%
普通VM中断。这里的普通VM中断，是指timer interrupt、I/O完成%
interrupt、network interrupt等guest OS正常运行时会看到的中断，%
它们不同于用于结束epoch的控制中断。%
\\对于primary和backup来说，它们从相同状态启动，并且各自的%
hypervisor都会将该特殊寄存器设置为相同的值，例如1000。%
在一个epoch内，primary收到的普通VM中断不会立即投递给%
primary VM，而是先由primary hypervisor记录并缓存起来；%
backup本机收到的普通VM中断也不会直接投递给backup VM，%
而是由backup hypervisor丢弃。%
primary VM和backup VM会执行对应的epoch，但backup至少落后primary%
一个epoch。当primary的epoch结束时，primary hypervisor把本epoch内%
缓存的普通VM中断列表发送给backup，并等待ack。backup执行到对应epoch%
边界后，等待并接收该中断列表，先向primary返回ack，再向backup VM投递%
同一批中断。primary收到ack后，才在自己的epoch边界向primary VM投递%
这批中断；随后双方分别进入下一个对应的epoch。}%

\subsection{FT协议}

\insertfiguretwo

对于VMware FT，我们使用deterministic replay来生成必要的日志条目，%
用来记录primary VM的执行过程。%
但我们不是把这些log entry写入磁盘，而是通过logging channel把%
它们发送给backup VM。%
backup VM会实时replay这些log entry，%
因此以与primary VM完全相同的方式执行。不过，为了确保真正实现fault %
tolerance，我们必须在logging channel上的log entry之%
外增加严格的FT协议。我们的基本要求如下：%

\begin{quote}
\noindent\textbf{Output Requirement:}如果backup VM在primary VM %
failure后take over，那么backup VM将继续以一种完全符合%
primary VM已经发送到外部世界的所有输出的方式执行。%
\end{quote}

\noindent{}注意，在发生failover之后(也就是backup VM在primary VM %
failure后take over之后)，由于执行过程中会发生许多%
non-deterministic event，backup VM很可能开始以一种与%
primary VM原本继续执行时非常不同的方式执行。%
不过，只要backup VM满足\textbf{Output Requirement}，%
在failover到backup VM的过程中就不会丢失任何外部可见的状态或数据，%
客户端也不会察觉服务中断或不一致。%

\textbf{Output Requirement}可以通过delay任何外部输出%
(通常是一个网络包)来保证，直到backup VM已经接收到所有%
必要信息，使其能够至少将执行重放到该输出操作发生的位置。%
一个必要条件是，backup VM必须已经接收到该输出操作之前生%
成的所有日志条目。这些日志条目将使它能够执行到最后一个日志%
条目所在的位置。然而，假设primary VM在执行完输出操作后立即发生%
故障，backup VM就必须知道自己必须继续replay，直到该输出操%
作的位置，并且\textbf{只能在那个位置“上线”}(也就是停止重放，并接%
管成为primary VM，见section 2.3)。如果备份虚拟机在输出操作之%
前的最后一个日志条目处就上线(go live)，那么某些non-deterministic event%
(例如传递给VM的timer interrupt)可能会在它执行该输出%
操作之前改变它的执行路径。
\\\emph{
译者注：具体来说，常见log entry包括：
\begin{enumerate}
  \item \textbf{input event log}：记录外部输入及其投递位置，例如incoming network packet、
  disk read completion、键盘输入、鼠标输入等。这类事件不仅表现为interrupt，
  还会把外部数据带入VM，因此必须记录数据内容和投递位置。
  \item \textbf{interrupt/event log}：记录不携带外部数据、但会改变guest执行路径的
  non-deterministic event，例如timer interrupt、disk write completion interrupt、
  以及其他设备completion interrupt。这类事件的重点是记录它在指令流中的精确投递位置。
  \item \textbf{non-deterministic instruction log}：记录某些非确定性指令或特殊读操作的结果，
  例如读取时间戳计数器、时钟、某些设备状态寄存器等。如果guest执行这类操作时得到的结果
  不能由初始状态和普通指令流确定，那么replay时也必须使用primary记录下来的结果。
  \item \textbf{output log entry}：记录external output operation在指令流中的发生位置，
  例如network packet send，必要时也包括某些对外可见的磁盘写操作。它的核心作用不是让
  backup在replay时再次对外输出，而是把output point归一化为log stream中的一个
  marker，使backup知道自己必须至少replay到这个位置，才能安全go live。
\end{enumerate}
}

在上述约束下，强制执行\textbf{Output Requirement}最简单的方法，%
是在每个output operation处创建一个特殊的log entry。于是，%
\textbf{Output Requirement}可以通过下面这条具体规则来强制执行：%

\begin{quote}
\noindent\textbf{Output Rule:}在backup VM收到并acknowledge与产生该%
output的操作相关联的log entry之前，primary VM不得向外部世界发送该输出%
(向外部世界发送输出包括网络packet输出和磁盘操作)。%
\end{quote}

如果backup VM已经收到所有log entry，包括产生输出的操作对应的log %
entry，那么backup VM就能够精确复现primary VM在该输出点处的状%
态；因此，如果primary VM死亡，backup VM将正确到达一个与该输出一致%
的状态。相反，如果backup VM在没有收到所有必要log entry的情况下%
take over，那么它的状态可能很快发散，导致与primary VM的输出不一致%
。\textbf{Output Rule}在某种程度上类似于[11]中描述的方法，%
其中“外部同步”的IO实际上可以被缓冲，只要它在下一次外部通信之前确实写入磁盘即%
可。%

注意，\textbf{Output Rule}并没有要求停止primary VM的执行。%
我们只需要delay发送输出，而VM本身可以继续执行\emph{(译者注：
这也是为什么这个系统能较少地影响机器性能的部分原因，不会像之前提%
到的epoch系统影响到primary正常运行)}。由于操作系统使用非阻塞的网%
络和磁盘输出，并通过异步interrupt表示完成，VM可以很容易继续执行，%
并且不一定会立即受到输出delay的影响。相比之下，以前的工作[3, 9]通常指%
出，primary VM在执行输出之前必须完全停止，直到backup VM %
acknowledge已收到来自primary VM的所有必要信息。\emph{
(译者注：比如前面的epoch，primary和backup都会停在执行完某个epoch
之后，然后这里会暂停primary VM，等到交互完成后双方才能继续执行。
因此这个系统虽然是“批处理”思想的系统，但也是性能损失较大的)}%

作为示例，我们在Figure 2中展示了一张说明FT协议要求的图。该图展示了%
primary VM和backup VM上事件的时间线。从primary VM线指向%
backup VM线的箭头表示log entry的传输，从backup VM线指向%
primary VM线的箭头表示acknowledgment。异步事件、%
输入和output operation的信息都必须作为log entry发送到%
backup VM，并得到acknowledgment。如图所示，%
到外部世界的输出会被delay，直到primary VM收到backup VM的%
acknowledgment，表明backup VM已经收到与该output %
operation相关联的log entry。%
只要遵循Output Rule，backup VM就能够以一种与primary VM最%
后一次输出一致的状态take over。%

我们无法保证在failover场景中所有输出都恰好产生一次。%
如果在primary VM打算发送输出时不使用带two-phase commit的%
transaction，backup VM就无法判断primary VM是在发送最后%
一个输出之前还是之后崩溃。幸运的是，network infrastructure(%
包括TCP的常见使用方式)本来就是为了处理丢失包和相同的(重复)包而设计的。%
注意，primary VM的传入包也可能在primary VM failure期间丢失%
，因此不会被投递到backup VM。不过，传入包可能因为许多与服务器%
failure无关的原因被丢弃，因此network infrastructure、%
操作系统和应用都被编写成能够补偿丢包。%
\\\emph{译者注：
\\这里并不能保证packet恰好发送一次。primary可能在最后一个packet%
真正发出之前或之后failure，因此客户端可能观察到丢包或一个包。%
backup决定上任后会切断log channel，然后重放log entry，%
对于发送网络报文的log entry直接忽略，因此这可能丢包，或者发出一个包；
当然重放发送网络报文的log entry时重新发包也是正确的实现，在这种实现中，%
真实发出一个包或者两个包都是有可能的。无论是哪种实现，都是正确的。%
\\\\Output Rule保证的是：只要某个输出已经对外可见，backup就拥有足够的%
log，可以在failover后到达与该输出一致的状态。TCP等网络协议本来就会%
处理丢包、重复和乱序，因此FT层不需要判断最后一个packet究竟是否发出。%
相比于本文主要描述的共享磁盘(它本质上也是一种输出设备)，%
这种磁盘I/O没有上述网络协议的语义，需要额外考虑和处理，相关处理见section 3.4。}%

\subsection{检测并响应failure}

如上所述，如果另一个VM看起来已经failure\emph{(primary%
和backup都需要感知对方，对面挂了就采取一定措施)}，primary VM和%
backup VM必须快速响应。如果backup VM failure，%
primary VM将go live，也就是离开recording mode(因此停止%
在logging channel上发送log entry)并开始正常执行%
\emph{(译者注：实际上这里说primary go live用词不妥，它的意思%
仅仅是退出recording模式，变成独立运行/不受保护的模式，%
并不代表之前就不是“live”的，它仅仅是说后续不再需要得到backup的ack确认后%
才能对外输出，而是放弃容错，确保至少能对外提供服务)}。%
如果primary VM failure，backup VM也应类似地go live，%
但这个过程稍微复杂一些。由于backup VM的执行存在滞后，%
它很可能已经接收并acknowledge了若干log entry，%
但由于尚未到达执行中的相应位置，这些log entry还没有被消费而被%
堆积。backup VM须继续执行并消费日志，直到最后一条log entry被重放%
\emph{(译者注：log entry类型众多，读者可以去2.2的译者注查阅)}，%
此时，backup VM将停止replay模式，并开始作为普通VM执行。%
本质上，backup VM已被提升为primary VM(但此时它缺少一个backup VM)。%
由于它不再是backup VM，新的primary VM现在会在guest OS执行%
output operation时向外部世界产生输出。在切换到正常模式期间，%
可能需要进行一些设备相关的操作，才能使这些输出正确发生。特别是对于网络，%
VMware FT会自动在网络上宣告新primary VM的MAC地址，%
使物理网络交换机知道新primary VM位于哪台服务器\emph{(%
译者注：这样后续来自客户端的网络报文才能继续被处理)}。此外，%
新提升的primary VM可能需要重新发出一些磁盘IO(如section 3.4所述)。%

尝试检测primary VM和backup VM failure的方法有很多。%
VMware FT使用运行fault-tolerant VM的服务器之间的UDP %
heartbeat来检测某台服务器是否可能已经崩溃。此外，%
VMware FT会监控从primary VM发送到backup VM的logging %
traffic，以及从backup VM发送到primary VM的%
acknowledgment。由于有定期的timer interrupt，%
对于正常运行的guest OS，logging traffic应当是规律的，%
且不会停止。因此，log entry或acknowledgment的流动停止可能表%
明某个VM发生failure。如果heartbeat或logging %
traffic停止超过某个特定超时时间(大约几秒)，就会宣布发生failure。%

然而，任何这样的failure detection方法都容易受到%
split-brain问题的影响。如果backup服务器停止接收来自primary%
服务器的heartbeat，这可能表示primary服务器已经failure，%
也可能只是表示仍在正常运行的服务器之间失去了所有网络连通性%
\emph{(译者注：也就是backup和primary发生分区，互相无法ping%
到对方，UDP的heartbeat心跳包也发不过去，都认为对方死亡)}。%
如果此时backup VM go live，而primary VM实际上仍在运行，%
那么很可能发生数据损坏，并给与VM通信的客户端带来问题。因此，%
当检测到failure时，我们必须确保primary VM和backup VM中只有%
一个go live。为了避免split-brain问题，我们利用存储VM %
virtual disk的shared storage。当primary VM或%
backup VM想要go live时，它会在shared storage上执行一个%
atomic的test-and-set操作。如果操作成功，该VM就被允许go live%
。如果操作失败，那么另一个VM一定已经go live，因此当前VM实际上会停止自身%
(“自杀”)。如果VM在尝试执行该atomic操作时无法访问shared %
storage，那么它只会等待，直到可以访问为止。注意，%
如果shared storage由于存储网络中的某种failure而不可访问，%
那么该VM很可能无论如何都无法做有用工作，因为virtual disk位于同一个%
shared storage上。因此，使用shared storage解决%
split-brain情况不会引入额外的不可用性。%
\\\emph{译者注：论文中的shared storage不是由guest OS%
直接访问的网络盘，也不是挂载在guest OS文件系统层中的%
对象。你可以把它简单理解为仅能在hypervisor层使用的%
共享块存储，它的接口很简单，只有Read(offset, length)%
和Write(offset, data)，看起来就是一块大的共享数组。%
hypervisor再把这个共享存储模拟成一个真正的磁盘设备，guest OS%
感知到的仍然是/dev/sda这样的本地块设备。%
\\\\当操作系统尝试提交IO请求的时候，会向ring放置IO Req，%
然后idx++，写入MMIO。VM会截获这个MMIO写入事件然后VM-Exit让%
VM获得控制权，此时它就会把guest OS的磁盘IO Req转换为共享存储%
读写。此外，test-and-set操作也是在hypervisor层进行的，%
guest OS无感知。}%

设计的最后一个方面是，一旦发生failure并且其中一个VM已经go live，%
VMware FT会通过在另一台host上启动新的backup VM来自动恢复冗余。%
虽然大多数先前工作没有覆盖这个过程，但它对于使fault-tolerant VM真%
正有用是根本性的，并且需要仔细设计。更多细节见section 3.1。
